\section*{Introduction}

\textbf{The Intelligent Investor} is not a guide to beating the market, but a timeless manual for cultivating the right mindset and principles for long-term investing. The book emphasizes \textbf{investing over speculation}, urging readers to adopt a disciplined, research-driven approach rather than trying to “time the market.”

\subsection*{Purpose and Philosophy}
The book is intended to guide investors in adopting and executing sound investment policies. Intelligent investors perform deep research on stocks and bonds to understand their behavior under varying conditions. Speculators, on the other hand, attempt to profit from short-term market movements, which Graham argues is contrary to common business sense and unsustainable.

\subsection*{Investor Types}
\begin{itemize}
    \item \textbf{Defensive Investor}: Focuses on avoiding serious mistakes or losses, and seeks freedom from effort, annoyance, and frequent decision-making.
    \item \textbf{Enterprising Investor}: Willing to devote time and care to selecting securities that are sound and more attractive than average.
\end{itemize}

\subsection*{Key Insights}
\begin{itemize}
    \item \(
        \mathrm{Growth} \neq \mathrm{Profit}
    \): Physical growth in an industry does not guarantee investor profits. For example, the airline industry has shown cumulative losses despite expansion.
    \item There is no reliable method to identify the most promising companies in a booming industry, such as picking IBM among computer firms in the 1970s.
    \item Investment principles remain constant over time, but their application must adapt to changing financial climates.
    \item Enthusiasm, while useful in other areas of life, often leads to disaster in investing.
    \item Cultivate the habit of measuring and quantifying value—consistently asking “How much?" and relating price to tangible asset value. This mindset, anchored in the principle of margin of safety, is essential for sound investing.

\end{itemize}

\subsection*{Main Objectives of the Book}
\begin{enumerate}
    \item Avoid irreversible losses.
    \item Maximize chances of achieving sustainable gains.
    \item Control self-defeating behavior.
\end{enumerate}

\subsection*{Commentary Highlights by Jason Zweig}
Zweig emphasizes that no book can truthfully tell you how to beat the market. He defines the intelligent investor as someone who is \textbf{patient, disciplined, eager to learn, emotionally resilient, and capable of independent thought}. This is more a matter of character than intellect.

\textit{Example:} Sir Isaac Newton, despite his brilliance, lost a fortune in the South Sea Company by letting crowd psychology override his judgment—demonstrating that emotional discipline is essential.

The market behaves like a pendulum, swinging between irrational exuberance and unjustifiable pessimism. Investors should welcome bear markets as opportunities to buy quality stocks at discounted prices—provided they have enough cash to meet their needs.